% GENERATED FILE. Edit canonical lessons, then run npm run generate:books.
% canonical-source-hash: fnv1a64:37d1e8b1d71077e9
% canonical-lessons: SA-C225-upatyaka, SA-C225-marubhumih, SA-C225-kananam, SA-C225-kunjah, SA-C225-sadvalam

\chapter{Valley, Desert, Woodland, Bower, Meadow}
\label{ch:sa-valley-desert-woodland-bower-meadow}
\hlchaptermodality{225}

\begin{chapteropening}
\textbf{By the end of this chapter:} \emph{I can say a valley, a desert, a woodland, a bower and a meadow.}

Complete the last lesson of chapter 225: I can say a valley, a desert, a woodland, a bower and a meadow.
\end{chapteropening}

\section[\texorpdfstring{\sk{उपत्यका} (upatyakā)}{उपत्यका (upatyakā)}]{\sk{उपत्यका} (upatyakā) — a valley}

\label{lesson:SA-C225-upatyaka}



\begin{quote}
Before the new one: say the Sanskrit for radiance, then the Sanskrit for the ground.
\end{quote}

\subsection*{You'll want to know: \sk{उपत्यका}}
\textbf{\sk{उपत्यका}} — \emph{upatyakā} — \textquotedblleft{}a valley\textquotedblright{}.

\begin{grammarlens}[title={What you've built}]
One more word for this chapter.
\end{grammarlens}

\subsection*{Your turn}
\begin{itemize}
\raggedright
  \item \emph{Say it:} \emph{upatyakā}
  \item \emph{Say it:} \emph{upatyakā}, once more
  \item \emph{Say it:} say \emph{upatyakā}
  \item \emph{Recall:} say the Sanskrit for radiance, then the Sanskrit for the ground, then say \emph{upatyakā} again
\end{itemize}

\begin{tcolorbox}[breakable,colback=teal!4,colframe=teal!35!black,title={Before you move on}]
What does \sk{उपत्यका} mean? (\textquotedblleft{}A valley\textquotedblright{}.) Say it once more.
\end{tcolorbox}
\section[\texorpdfstring{\sk{मरुभूमिः} (marubhūmiḥ)}{मरुभूमिः (marubhūmiḥ)}]{\sk{मरुभूमिः} (marubhūmiḥ) — a desert}

\label{lesson:SA-C225-marubhumih}



\begin{quote}
Before the new one: say the Sanskrit for the ground, then the Sanskrit for a valley.
\end{quote}

\subsection*{You'll want to know: \sk{मरुभूमिः}}
\textbf{\sk{मरुभूमिः}} — \emph{marubhūmiḥ} — \textquotedblleft{}a desert\textquotedblright{}.

\begin{grammarlens}[title={What you've built}]
One more word for this chapter.
\end{grammarlens}

\subsection*{Your turn}
\begin{itemize}
\raggedright
  \item \emph{Say it:} \emph{marubhūmiḥ}
  \item \emph{Say it:} \emph{marubhūmiḥ}, once more
  \item \emph{Say it:} say \emph{marubhūmiḥ}
  \item \emph{Recall:} say the Sanskrit for the ground, then the Sanskrit for a valley, then say \emph{marubhūmiḥ} again
\end{itemize}

\begin{tcolorbox}[breakable,colback=teal!4,colframe=teal!35!black,title={Before you move on}]
What does \sk{मरुभूमिः} mean? (\textquotedblleft{}A desert\textquotedblright{}.) Say it once more.
\end{tcolorbox}
\section[\texorpdfstring{\sk{काननम्} (kānanam)}{काननम् (kānanam)}]{\sk{काननम्} (kānanam) — a woodland}

\label{lesson:SA-C225-kananam}



\begin{quote}
Before the new one: say the Sanskrit for a valley, then the Sanskrit for a desert.
\end{quote}

\subsection*{You'll want to know: \sk{काननम्}}
\textbf{\sk{काननम्}} — \emph{kānanam} — \textquotedblleft{}a woodland\textquotedblright{}.

\begin{grammarlens}[title={What you've built}]
One more word for this chapter.
\end{grammarlens}

\subsection*{Your turn}
\begin{itemize}
\raggedright
  \item \emph{Say it:} \emph{kānanam}
  \item \emph{Say it:} \emph{kānanam}, once more
  \item \emph{Say it:} say \emph{kānanam}
  \item \emph{Recall:} say the Sanskrit for a valley, then the Sanskrit for a desert, then say \emph{kānanam} again
\end{itemize}

\begin{tcolorbox}[breakable,colback=teal!4,colframe=teal!35!black,title={Before you move on}]
What does \sk{काननम्} mean? (\textquotedblleft{}A woodland\textquotedblright{}.) Say it once more.
\end{tcolorbox}
\section[\texorpdfstring{\sk{कुञ्जः} (kuñjaḥ)}{कुञ्जः (kuñjaḥ)}]{\sk{कुञ्जः} (kuñjaḥ) — a bower, a thicket}

\label{lesson:SA-C225-kunjah}



\begin{quote}
Before the new one: say the Sanskrit for a desert, then the Sanskrit for a woodland.
\end{quote}

\subsection*{You'll want to know: \sk{कुञ्जः}}
\textbf{\sk{कुञ्जः}} — \emph{kuñjaḥ} — \textquotedblleft{}a bower, a thicket\textquotedblright{}.

\begin{grammarlens}[title={What you've built}]
One more word for this chapter.
\end{grammarlens}

\subsection*{Your turn}
\begin{itemize}
\raggedright
  \item \emph{Say it:} \emph{kuñjaḥ}
  \item \emph{Say it:} \emph{kuñjaḥ}, once more
  \item \emph{Say it:} say \emph{kuñjaḥ}
  \item \emph{Recall:} say the Sanskrit for a desert, then the Sanskrit for a woodland, then say \emph{kuñjaḥ} again
\end{itemize}

\begin{tcolorbox}[breakable,colback=teal!4,colframe=teal!35!black,title={Before you move on}]
What does \sk{कुञ्जः} mean? (\textquotedblleft{}A bower, a thicket\textquotedblright{}.) Say it once more.
\end{tcolorbox}
\section[\texorpdfstring{\sk{शाद्वलम्} (śādvalam)}{शाद्वलम् (śādvalam)}]{\sk{शाद्वलम्} (śādvalam) — a meadow}

\label{lesson:SA-C225-sadvalam}



\begin{quote}
Before the new one: say the Sanskrit for a woodland, then the Sanskrit for a bower.
\end{quote}

\subsection*{You'll want to know: \sk{शाद्वलम्}}
\textbf{\sk{शाद्वलम्}} — \emph{śādvalam} — \textquotedblleft{}a meadow\textquotedblright{}.

\begin{grammarlens}[title={What you've built}]
One more word for this chapter.
\end{grammarlens}

\subsection*{Your turn}
\begin{itemize}
\raggedright
  \item \emph{Say it:} \emph{śādvalam}
  \item \emph{Say it:} \emph{śādvalam}, once more
  \item \emph{Say it:} say \emph{śādvalam}
  \item \emph{Recall:} say the Sanskrit for a woodland, then the Sanskrit for a bower, then say \emph{śādvalam} again
\end{itemize}

\begin{tcolorbox}[breakable,colback=teal!4,colframe=teal!35!black,title={Before you move on}]
What does \sk{शाद्वलम्} mean? (\textquotedblleft{}A meadow\textquotedblright{}.) Say it once more.
\end{tcolorbox}
